The reference implementation is an open-source library cited as
\texttt{{[}anonymised\ for\ blind\ review{]}}. Approximately 12,000
lines of Python implement the runtime (the SMT prover, capability
issuance and the gate, provenance receipts and chain verification, the
modelled-language registry described below, the audit chain, the
policy-DSL gateway, remote-key signing, and the conformance
re-verification layer described below); an additional \textasciitilde430
lines of Lean 4 contain the formal development. The implementation
depends on Z3 (Python bindings) and the audited \texttt{cryptography}
library for Ed25519. No new cryptographic primitives are introduced.

Several engineering decisions warrant comment.

\textbf{Modelled-language registry.} The set of constructs the system
reasons about --- capability constraint kinds, JSON-Schema keywords, the
URL and SQL grammar/policy surfaces, and the provenance chain-envelope
fields --- is enumerated in a single executable registry. Both the
runtime and the test suite derive their allow-lists from it, and a
continuous-integration drift test fails if any consumer's modelled set
diverges from the registry. This makes fail-closed behaviour a
\emph{structural} property rather than a coding convention: a construct
absent from the registry hits its dimension's conservative verdict by
construction. An unmodelled schema keyword, policy key, URL key, or SQL
construct downgrades a would-be \texttt{PROVEN} to \texttt{UNDECIDED};
an unknown capability constraint kind is rejected at mint; an unknown
chain-envelope field is rejected at verification. The registry is the
structural mechanism behind Theorem 2's ``modelled constraint kinds''
qualifier: the gate cannot silently enforce, or fail to enforce, a kind
the model does not name.

\textbf{Canonical JSON.} All hashes and signature inputs use a
canonical-JSON serialisation --- sorted keys, no whitespace, UTF-8
encoded, \texttt{ensure\_ascii=False} --- so that the same logical
object always produces the same bytes. This avoids signature-mismatch
failures that bedevil deployments where serialisation order varies
across implementations. The canonical form is a restricted JCS subset:
integers only (floats are rejected rather than serialised, because
cross-implementation float canonicalisation is the one genuinely hard
part of JCS), bounded to the JavaScript safe-integer range so that every
value round-trips byte-identically across all five reference
implementations (Python, TypeScript, Go, Rust, C\#). A conformance
harness checks byte-identity on emit and verify across the five; the
provenance chain envelope is the minimal
\texttt{\{receipt\_hash,\ jws\}} pair, with the signed payload carried
inside the JWS rather than duplicated --- and trusted --- in the
envelope.

\textbf{Fail-closed defaults.} The gate raises rather than returning
ALLOW on any internal error. The prover returns \texttt{UNDECIDED}
(treated as a failure mode by downstream consumers) rather than
\texttt{PROVEN} on solver timeout. The principle is applied uniformly on
every verification boundary: unmodelled inputs map to the conservative
verdict above, a token that cites a parent it cannot resolve is denied
rather than trusted, and verification paths reject non-conforming input
rather than warning-and-accepting it. We observed in early integration
testing that these defaults catch a class of misconfiguration bugs --- a
missing pubkey in the trusted-issuer map, a corrupted token on disk ---
that fail-open defaults would silently mask.

\textbf{Out-of-process gate.} Although the gate can run as a Python
library function in the same address space as the agent runtime, the
recommended deployment runs it as a separate process behind a Unix
socket or local HTTPS endpoint, with the issuer key held by an HSM. This
separation provides defence-in-depth: an LLM-driven RCE in the agent
runtime cannot directly mint or modify tokens.

\textbf{Content-addressed everything.} Tokens, proof artifacts,
audit-chain leaves, and feed entries (where used in a larger deployment)
are all content-addressed, allowing cross-reference by hash without
needing a central registry.

\textbf{Conformance re-verification with selective disclosure.} The gate
records \texttt{args\_hash\ =\ SHA-256(canonical(args))} in each
receipt, so a third-party verifier can confirm signatures, cited schema
and proof hashes, and chain integrity offline --- but cannot confirm
that \emph{this specific call's arguments} satisfied the policy without
the operator disclosing raw arguments. The conformance re-verification
layer closes this gap for auditable subsets. Call arguments are
committed as a Merkle tree over their top-level fields; leaves hash the
field name and the JCS-canonical encoding of the value; ordering follows
UTF-16 code-unit order, matching the canonicalisation discipline proven
cross-language for receipt bodies. The operator discloses only
constraint-relevant fields (with Merkle paths); the verifier recomputes
and re-derives per-constraint SATISFIED / VIOLATED / UNKNOWN verdicts
independently of the operator's runtime, failing closed on any
undisclosed field. All hashing is SHA-256, chosen deliberately because
it is the well-trodden path in zk proof systems: the scheme ports to a
zkVM (Layer 2) without exotic primitives.

\textbf{Remote-key signing.} In production the issuer key need not live
in a local HSM: the signer abstraction routes Ed25519 operations to AWS
KMS, Azure Key Vault, or HashiCorp Vault Transit, so the key material
never leaves the remote service. The token, proof, and audit artefacts
are unchanged; only the signing substrate differs. The gate is agnostic
to which signer minted a token, provided its \texttt{key\_id} is in the
trusted-issuer map.

The codebase is released under a permissive open-source licence
(Apache-2.0); the Lean development will be released alongside
camera-ready acceptance.
